% GENERATED FILE. Edit canonical lessons, then run npm run generate:books.
% canonical-source-hash: fnv1a64:be79521f9e955b51
% canonical-lessons: IT-C348-vantaggio, IT-C348-svantaggio, IT-C348-aspetto, IT-C348-critica, IT-C348-ribadire

\chapter{Pro e contro}
\label{ch:it-pro-e-contro}
\hlchaptermodality{348}

\begin{chapteropening}
\textbf{By the end of this chapter:} \emph{I can weigh up pros and cons: name an advantage, a disadvantage and an aspect, voice a criticism and repeat my point.}

Complete the last lesson of chapter 348: I can weigh up pros and cons: name an advantage, a disadvantage and an aspect, voice a criticism and repeat my point.
\end{chapteropening}

\section[il vantaggio]{il vantaggio — the advantage}

\label{lesson:IT-C348-vantaggio}



\begin{quote}
Before the new one: say the Italian for objective, then the Italian for subjective.
\end{quote}

\subsection*{You'll want to know: il vantaggio}
\textbf{il vantaggio} — \textquotedblleft{}the advantage\textquotedblright{}.

\begin{quote}
Abitare in centro ha molti vantaggi.
\end{quote}

\begin{quote}
\emph{Living in the centre has many advantages.}
\end{quote}

\emph{Essere in vantaggio} — to be ahead, in a match too.

\begin{cousinweb}
\emph{vantaggio} is usually traced to Old French \emph{avantage}, built on \emph{avant} \textquotedblleft{}before, ahead\textquotedblright{}, from Late Latin \emph{abante}, made from \emph{ab} and \emph{ante} \textquotedblleft{}before\textquotedblright{} (Lewis \& Short): being a step ahead.
\end{cousinweb}

\begin{grammarlens}[title={What you've built}]
One more word for telling, explaining and describing at B1.
\end{grammarlens}

\subsection*{Your turn}
\begin{itemize}
\raggedright
  \item \emph{Say it:} \emph{il vantaggio}
  \item \emph{Say it:} \emph{il vantaggio}, once more
  \item \emph{Say it:} a sentence of your own with \emph{il vantaggio}
  \item \emph{Recall:} say the Italian for objective, then the Italian for subjective, then say \emph{il vantaggio} again
\end{itemize}

\begin{tcolorbox}[breakable,colback=teal!4,colframe=teal!35!black,title={Before you move on}]
What does il vantaggio mean? (\textquotedblleft{}The advantage\textquotedblright{}.) Say it once more.
\end{tcolorbox}
\section[lo svantaggio]{lo svantaggio — the disadvantage}

\label{lesson:IT-C348-svantaggio}



\begin{quote}
Before the new one: say the Italian for subjective, then the Italian for the advantage.
\end{quote}

\subsection*{You'll want to know: lo svantaggio}
\textbf{lo svantaggio} — \textquotedblleft{}the disadvantage\textquotedblright{}.

\begin{quote}
Lo svantaggio è che la casa è lontana dalla stazione.
\end{quote}

\begin{quote}
\emph{The disadvantage is that the house is far from the station.}
\end{quote}

\emph{S-} in front reverses the meaning, as in \emph{comodo} and \emph{scomodo}.

\begin{cousinweb}
\emph{svantaggio} is built from \emph{s-} \textquotedblleft{}un-\textquotedblright{} and \emph{vantaggio}, which is usually traced to Old French \emph{avantage}, from Late Latin \emph{abante} \textquotedblleft{}ahead\textquotedblright{}, made from \emph{ab} and \emph{ante} (Lewis \& Short).
\end{cousinweb}

\begin{grammarlens}[title={What you've built}]
One more word for telling, explaining and describing at B1.
\end{grammarlens}

\subsection*{Your turn}
\begin{itemize}
\raggedright
  \item \emph{Say it:} \emph{lo svantaggio}
  \item \emph{Say it:} \emph{lo svantaggio}, once more
  \item \emph{Say it:} a sentence of your own with \emph{lo svantaggio}
  \item \emph{Recall:} say the Italian for subjective, then the Italian for the advantage, then say \emph{lo svantaggio} again
\end{itemize}

\begin{tcolorbox}[breakable,colback=teal!4,colframe=teal!35!black,title={Before you move on}]
What does lo svantaggio mean? (\textquotedblleft{}The disadvantage\textquotedblright{}.) Say it once more.
\end{tcolorbox}
\section[l'aspetto]{l'aspetto — the aspect, the side (of a question); the look, the appearance}

\label{lesson:IT-C348-aspetto}



\begin{quote}
Before the new one: say the Italian for the advantage, then the Italian for the disadvantage.
\end{quote}

\subsection*{You'll want to know: l'aspetto}
\textbf{l'aspetto} — \textquotedblleft{}the aspect, the side (of a question); the look, the appearance\textquotedblright{}.

\begin{quote}
C'è un aspetto del problema che non abbiamo considerato.
\end{quote}

\begin{quote}
\emph{There is one side of the problem that we have not considered.}
\end{quote}

Also how something looks: \emph{ha un bell'aspetto}. Do not mix it up with \emph{aspetto} \textquotedblleft{}I wait\textquotedblright{}.

\begin{cousinweb}
\emph{aspetto} is a learned borrowing from Latin \emph{aspectus} \textquotedblleft{}a look, an appearance\textquotedblright{} (Lewis \& Short), from \emph{aspicere} \textquotedblleft{}to behold\textquotedblright{}, built on \emph{specere} \textquotedblleft{}to look\textquotedblright{}. The verb \emph{aspettare} \textquotedblleft{}to wait\textquotedblright{} comes from the same Latin family.
\end{cousinweb}

\begin{grammarlens}[title={What you've built}]
One more word for telling, explaining and describing at B1.
\end{grammarlens}

\subsection*{Your turn}
\begin{itemize}
\raggedright
  \item \emph{Say it:} \emph{l'aspetto}
  \item \emph{Say it:} \emph{l'aspetto}, once more
  \item \emph{Say it:} a sentence of your own with \emph{l'aspetto}
  \item \emph{Recall:} say the Italian for the advantage, then the Italian for the disadvantage, then say \emph{l'aspetto} again
\end{itemize}

\begin{tcolorbox}[breakable,colback=teal!4,colframe=teal!35!black,title={Before you move on}]
What does l'aspetto mean? (\textquotedblleft{}The aspect, the side (of a question); the look, the appearance\textquotedblright{}.) Say it once more.
\end{tcolorbox}
\section[la critica]{la critica — the criticism; the critics}

\label{lesson:IT-C348-critica}



\begin{quote}
Before the new one: say the Italian for the disadvantage, then the Italian for the aspect.
\end{quote}

\subsection*{You'll want to know: la critica}
\textbf{la critica} — \textquotedblleft{}the criticism; the critics\textquotedblright{}.

\begin{quote}
Accetto la critica, ma il lavoro è finito.
\end{quote}

\begin{quote}
\emph{I accept the criticism, but the work is finished.}
\end{quote}

\emph{Fare una critica} — to make a criticism. \emph{La critica} can also mean the critics as a group.

\begin{cousinweb}
\emph{critica} is built on \emph{critico}, a learned borrowing from Latin \emph{criticus} \textquotedblleft{}able to judge\textquotedblright{} (Lewis \& Short), from Greek \emph{kritikós}, from \emph{krínein} \textquotedblleft{}to separate, to decide\textquotedblright{}.
\end{cousinweb}

\begin{grammarlens}[title={What you've built}]
One more word for telling, explaining and describing at B1.
\end{grammarlens}

\subsection*{Your turn}
\begin{itemize}
\raggedright
  \item \emph{Say it:} \emph{la critica}
  \item \emph{Say it:} \emph{la critica}, once more
  \item \emph{Say it:} a sentence of your own with \emph{la critica}
  \item \emph{Recall:} say the Italian for the disadvantage, then the Italian for the aspect, then say \emph{la critica} again
\end{itemize}

\begin{tcolorbox}[breakable,colback=teal!4,colframe=teal!35!black,title={Before you move on}]
What does la critica mean? (\textquotedblleft{}The criticism; the critics\textquotedblright{}.) Say it once more.
\end{tcolorbox}
\section[ribadire]{ribadire — to reiterate, to restate firmly}

\label{lesson:IT-C348-ribadire}



\begin{quote}
Before the new one: say the Italian for the aspect, then the Italian for the criticism.
\end{quote}

\subsection*{You'll want to know: ribadire}
\textbf{ribadire} — \textquotedblleft{}to reiterate, to restate firmly\textquotedblright{}.

\begin{quote}
Ribadisco quello che ho detto: io non vengo.
\end{quote}

\begin{quote}
\emph{I repeat what I said: I am not coming.}
\end{quote}

Conjugated with \emph{-isc-}: \emph{ribadisco}. The first sense is to clinch a nail by bending its point.

\begin{cousinweb}
The origin of \emph{ribadire} is not settled; EtymDB records the word with no source (EtymDB). Its first sense was \textquotedblleft{}to clinch a nail by hammering its point flat\textquotedblright{}, and from there it came to mean to hammer a point home with words.
\end{cousinweb}

\begin{grammarlens}[title={What you've built}]
One more word for telling, explaining and describing at B1.
\end{grammarlens}

\subsection*{Your turn}
\begin{itemize}
\raggedright
  \item \emph{Say it:} \emph{ribadire}
  \item \emph{Say it:} \emph{ribadire}, once more
  \item \emph{Say it:} a sentence of your own with \emph{ribadire}
  \item \emph{Recall:} say the Italian for the aspect, then the Italian for the criticism, then say \emph{ribadire} again
\end{itemize}

\begin{tcolorbox}[breakable,colback=teal!4,colframe=teal!35!black,title={Before you move on}]
What does ribadire mean? (\textquotedblleft{}To reiterate, to restate firmly\textquotedblright{}.) Say it once more.
\end{tcolorbox}
